\section{Time Reversibility}

\begin{definition}[Reversibility]
    A chain with transition matrix \(P\) and invariant distribution \(\pi\) is said to be \term{reversible}, if
    \[
        P\pqty{x, y} = \hat{P}\pqty{x, y},
    \]
    where \(\hat{P}\) is defined by
    \[
        \hat{P}\pqty{x, y} = \frac{\pi\pqty{y}}{\pi\pqty{x}} P\pqty{y, x},
    \]
    for all \(x, y \in I\).
\end{definition}

Note in particular that \(\hat{P}\) is always a transition matrix, with
\[
    \sum_{y \in I} \hat{P}\pqty{x, y} = \pqty{\sum_y \pi\pqty{y} P\pqty{y, x}} \frac{1}{\pi\pqty{x}} = 1
\]
for all \(x \in I\).

\begin{theorem}[Reversibility and Irreducible Chains]
    Suppose \(X \sim \Markov\pqty{\pi, P}\) is irreducible and \(\pi\) is invariant for \(P\). For a fixed \(N \geq 1\), define \(Y_k = X_{N - k}\). Then the process \(Y = \pqty{Y_k}_{0 \leq k \leq N}\) satisfies \(Y \sim \Markov\pqty{\pi, \hat{P}}\), and is also irreducible.
\end{theorem}

\begin{proof}
    Since \(\pi\) is invariant and \(X_0 \sim \pi\), we have \(X_N = Y_0 \sim \pi\) for the initial distribution.

    To see it is Markov, we have
    \begin{align*}
        \Prob\pqty{Y_0 = y_0, \ldots, Y_N = y_N} &= \Prob\pqty{X_0 = y_N, \ldots, X_N = y_0} \\
        &= \pi\pqty{y_N} P\pqty{y_N, y_{N - 1}} \cdots P\pqty{y_1, y_0} \\
        &= \pi\pqty{y_0} \hat{P}\pqty{y_0, y_1} \cdots \hat{P}\pqty{y_{N - 1}, y_N} \\
        &= \Prob\pqty{Y_0 = y_0} \hat{P}\pqty{y_0, y_1} \cdots \hat{P}\pqty{y_{N - 1}, y_N},
    \end{align*}
    so \(Y \sim \Markov\pqty{\pi, \hat{P}}\).

    Since all \(\pi_i > 0\), we have \(\hat{P}\pqty{x, y} > 0\) if and only if \(P\pqty{y, x} > 0\). Hence \(X\) is irreducible implies \(Y\) is irreducible.
\end{proof}

\begin{lemma}[Reversibility Implies Invariance]
    If a measure \(\lambda\) satisfies
    \[
        \lambda_i P\pqty{i, j} = \lambda_j P\pqty{j, i}
    \]
    for all \(i, j \in I\), then \(\lambda\) is invariant for \(P\).
\end{lemma}

\begin{proof}
    Assuming this condition, and summing over all \(i\), we have
    \[
        \sum_{i \in I} \lambda_i P\pqty{i, j} = \lambda_j \sum_{i \in I} P\pqty{j, i} = \lambda_j,
    \]
    for any \(j \in I\), so \(\lambda\) is invariant for \(P\).
\end{proof}

\begin{examples}
    \begin{enumerate}
        \item Suppose \(X\) is a random walk on the circle \(\ZZ_n = \set{0, 1, \ldots, n - 1}\) with \(P\pqty{i, i + 1} = \flatfrac{2}{3}\), and \(P\pqty{i, i - 1} = \flatfrac{1}{3}\) (all modulo \(n\)). Then
        \[
            \pi_i = \pi_{i + 1} \cdot \frac{1}{3} + \pi_{i - 1} \cdot \frac{2}{3},
        \]
        so \(\pi_i = \flatfrac{1}{n}\) is an invariant distribution, and since this chain is irreducible, it is unique. Moreover, it satisfies
        \[
            \pi_i \cdot \frac{2}{3} \neq \pi_{i + 1} \cdot \frac{1}{3},
        \]
        so the chain is not reversible.

        \item Suppose \(X\) is a random walk on \(\ZZ_n\), but this time with transition probabilities
        \begin{align*}
            P\pqty{i, i + 1} &= \frac{2}{3}, &  P\pqty{i, i - 1} &= \frac{1}{3} && 1 \leq i \leq n - 2, \\
            P\pqty{0, 0} &= \frac{1}{3}, & P\pqty{0, 1} &= \frac{2}{3}, \\
            P\pqty{n - 1, n - 1} &= \frac{2}{3}, & P\pqty{n - 1, n - 2} &= \frac{1}{3}. \\
        \end{align*}

        For this chain, it is easy to check that \(\lambda_i = 2^i\) satisfies the balance conditions
        \[
            \lambda_i P\pqty{i, j} = \lambda_{j} P\pqty{j, i},
        \]
        so this chain is reversible.

        \item Suppose \(G = \pqty{V, E}\) is a finite, simply connected, undirected graph, and let \(\deg x\) be the degree of vertex \(x\).  Consider the transition matrix \(P\) defined by
        \[
            P\pqty{x, y} = \begin{cases}
                \frac{1}{\deg x}, & \qif \set{x, y} \in E, \\
                0, & \quad\text{otherwise}.
            \end{cases}
        \]

        Then the distribution
        \[
            \pi_x = \frac{\deg x}{2 \abs{E}}
        \]
        is invariant for \(P\), and satisfies the detailed balance conditions. If \(\pqty{x, y} \in E\), we have
        \[
            \pi_x P\pqty{x, y} = \frac{\deg \pqty{x}}{2 \abs{E}} \frac{1}{\deg x} = \frac{1}{2 \abs{E}},
        \]
        and similarly,
        \[
            \pi_y P\pqty{y, x} = \frac{\deg \pqty{y}}{2 \abs{E}} \frac{1}{\deg y} = \frac{1}{2 \abs{E}}.
        \]

        Since the chain is irreducible, this is its unique invariant distribution, and the chain is reversible.
    \end{enumerate}
\end{examples}